\documentclass[12pt,a4paper]{article}

% Paquetes
\usepackage[utf8]{inputenc}
\usepackage[spanish]{babel}
\usepackage{geometry}
\usepackage{graphicx}
\usepackage{listings}
\usepackage{xcolor}
\usepackage{hyperref}
\usepackage{booktabs}
\usepackage{longtable}
\usepackage{array}
\usepackage{fancyhdr}
\usepackage{tcolorbox}

% Configuración de página
\geometry{
    left=2.5cm,
    right=2.5cm,
    top=3cm,
    bottom=3cm
}

% Configuración de código Python
\lstset{
    language=Python,
    basicstyle=\ttfamily\small,
    keywordstyle=\color{blue}\bfseries,
    commentstyle=\color{gray}\itshape,
    stringstyle=\color{red},
    showstringspaces=false,
    numbers=left,
    numberstyle=\tiny\color{gray},
    stepnumber=1,
    numbersep=5pt,
    backgroundcolor=\color{white},
    frame=single,
    rulecolor=\color{black},
    tabsize=4,
    breaklines=true,
    breakatwhitespace=false,
    captionpos=b,
    escapeinside={\%*}{*)},
    morekeywords={True,False,None}
}

% Configuración de hipervínculos
\hypersetup{
    colorlinks=true,
    linkcolor=blue,
    filecolor=magenta,
    urlcolor=cyan,
    pdftitle={Apéndice Técnico - Automatización HYSYS con Python},
    pdfauthor={Debugging Sistemático},
    pdfsubject={Referencia API HYSYS},
    pdfkeywords={HYSYS, Python, COM, Automatización, pywin32}
}

% Configuración de encabezado y pie de página
\pagestyle{fancy}
\fancyhf{}
\rhead{Apéndice Técnico - HYSYS API}
\lhead{Debugging Sistemático}
\rfoot{Página \thepage}

% Cajas de advertencia
\newtcolorbox{warningbox}{
    colback=red!5!white,
    colframe=red!75!black,
    title=\textbf{IMPORTANTE}
}

\newtcolorbox{successbox}{
    colback=green!5!white,
    colframe=green!75!black,
    title=\textbf{SOLUCIÓN}
}

\newtcolorbox{infobox}{
    colback=blue!5!white,
    colframe=blue!75!black,
    title=\textbf{HALLAZGO}
}

% Inicio del documento
\begin{document}

% Portada
\begin{titlepage}
    \centering
    \vspace*{2cm}

    {\Huge\bfseries APÉNDICE TÉCNICO\par}
    \vspace{1cm}
    {\LARGE Automatización de ASPEN HYSYS con Python\par}
    \vspace{0.5cm}
    {\Large Patrones de API Descubiertos mediante Debugging Sistemático\par}

    \vspace{3cm}

    \begin{tcolorbox}[colback=gray!10,colframe=gray!50,width=0.8\textwidth]
        \centering
        \textbf{Derivado de:} Sesión de debugging - Prácticas 1-8\\
        \textbf{Fecha:} 2025-01-20\\
        \textbf{Entorno:} HYSYS (compatible SRK), Python 64-bit, pywin32\\
        \textbf{Versión:} 1.0
    \end{tcolorbox}

    \vfill

    {\large Mantenedor: Claude (Anthropic)\par}
    {\large Estado: Validado mediante 8 prácticas funcionales\par}

    \vspace{1cm}
\end{titlepage}

% Tabla de contenidos
\tableofcontents
\newpage

% Sección 1: Resumen Ejecutivo
\section{Resumen Ejecutivo}

Durante la implementación de 8 prácticas de automatización de ASPEN HYSYS con Python, se encontraron múltiples incompatibilidades entre la documentación estándar y la API real expuesta por la versión instalada. Este documento registra los patrones correctos descubiertos mediante debugging sistemático.

\subsection{Contexto}

\begin{itemize}
    \item \textbf{Problema inicial:} Error COM -2147221021 al intentar agregar componentes
    \item \textbf{Duración:} Sesión extensa de debugging (iteración sistemática)
    \item \textbf{Resultado:} 15 archivos Python actualizados, 8 prácticas funcionales
    \item \textbf{Impacto:} Documentación de patrones reutilizables para implementaciones futuras
\end{itemize}

% Sección 2: Arquitectura de la API COM
\section{Arquitectura de la API COM de HYSYS}

\subsection{Jerarquía de Objetos Correcta}

La jerarquía real de objetos en la API COM de HYSYS difiere significativamente de la documentación estándar:

\begin{verbatim}
Application (HYSYS.Application)
    └── SimulationCases (colección)
        └── SimulationCase (objeto)
            ├── BasisManager ✓ (NO usar case.FluidPackage)
            │   ├── ComponentLists (colección)
            │   │   └── ComponentList (objeto)
            │   │       └── Components (colección)
            │   │           └── Component (objeto individual)
            │   │
            │   └── FluidPackages (colección)
            │       └── FluidPackage (objeto)
            │           └── PropertyPackageName (propiedad)
            │
            └── Flowsheet (objeto)
                ├── MaterialStreams (colección)
                └── Operations (colección)
\end{verbatim}

\subsection{Patrón Incorrecto vs Correcto}

\begin{warningbox}
\textbf{Patrón INCORRECTO} (documentación estándar):
\begin{lstlisting}
case = hysys.ActiveDocument  # Devuelve None
fluid_pkg = case.FluidPackage  # No existe
components = fluid_pkg.Components
components.Add("Water")  # Nombre incorrecto
\end{lstlisting}
\end{warningbox}

\begin{successbox}
\textbf{Patrón CORRECTO} (implementación real):
\begin{lstlisting}
case = hysys.SimulationCases.Add()
basis_manager = case.BasisManager
comp_lists = basis_manager.ComponentLists
comp_list = comp_lists.Add()
components = comp_list.Components
components.Add("H2O")  # Nombre correcto
\end{lstlisting}
\end{successbox}

% Sección 3: Patrones Específicos
\section{Patrones Específicos Descubiertos}

\subsection{Inicialización de Caso de Simulación}

\begin{lstlisting}
import win32com.client as win32

hysys = win32.Dispatch('HYSYS.Application')
hysys.Visible = True

# CORRECTO: Crear nuevo caso
case = hysys.SimulationCases.Add()

# INCORRECTO: ActiveDocument
# case = hysys.ActiveDocument  # Devuelve None
\end{lstlisting}

\begin{infobox}
\textbf{Hallazgo:} \texttt{ActiveDocument} requiere un caso previamente abierto. Para automatización completa, siempre usar \texttt{SimulationCases.Add()}.
\end{infobox}

\subsection{Agregación de Componentes}

El proceso de agregación de componentes requiere seguir estrictamente la jerarquía:

\begin{lstlisting}
# 1. Acceder al BasisManager
basis_manager = case.BasisManager

# 2. Crear una ComponentList
comp_lists = basis_manager.ComponentLists
comp_list = comp_lists.Add()

# 3. Acceder a Components DENTRO de la lista
components = comp_list.Components

# 4. Agregar componentes por nombre
components.Add("Methanol")
components.Add("H2O")  # NO "Water"
components.Add("Ethanol")

# 5. Verificar
print(f"Componentes agregados: {components.Count}")
\end{lstlisting}

\subsubsection{Tabla de Nombres de Componentes}

\begin{table}[h]
\centering
\begin{tabular}{@{}lll@{}}
\toprule
\textbf{Nombre Común} & \textbf{Nombre HYSYS Correcto} & \textbf{Alternativas (Fallan)} \\ \midrule
Water/Agua & \texttt{H2O} & Water, WATER, water \\
Methanol & \texttt{Methanol} & - \\
Ethanol & \texttt{Ethanol} & - \\
Ethyl acetate & \texttt{Ethyl acetate} & - \\
Glicerol & \texttt{Glycerol} & (no siempre disponible) \\ \bottomrule
\end{tabular}
\caption{Nombres correctos de componentes en HYSYS}
\label{tab:component_names}
\end{table}

\begin{warningbox}
\textbf{CRÍTICO:} El nombre \texttt{"Water"} NO está disponible en la base de datos. Siempre usar \texttt{"H2O"}.
\end{warningbox}

\subsection{Configuración de Paquete Termodinámico}

Los nombres de paquetes termodinámicos varían entre versiones de HYSYS. Implementar fallback:

\begin{lstlisting}
# Crear FluidPackage
fluid_pkg = basis_manager.FluidPackages.Add()

# Probar multiples nombres (varia entre versiones)
nombres_pkg = ["NRTL", "Peng-Robinson", "PR", "SRK",
               "Soave-Redlich-Kwong"]

for nombre in nombres_pkg:
    try:
        fluid_pkg.PropertyPackageName = nombre
        # Verificar que se asigno correctamente
        if fluid_pkg.PropertyPackageName == nombre:
            print(f"Modelo asignado: {nombre}")
            break
    except:
        continue  # Intentar siguiente nombre
\end{lstlisting}

\begin{infobox}
\textbf{Hallazgos:}
\begin{itemize}
    \item \texttt{SRK} funcionó en la versión de prueba
    \item \texttt{"Peng-Robinson"} (con guion) causó error -2147024809
    \item El nombre exacto varía entre versiones
    \item Siempre implementar fallback con múltiples variantes
\end{itemize}
\end{infobox}

\subsection{Extracción de Propiedades COM}

\subsubsection{Problema: Objetos CDispatch}

Las propiedades termodinámicas NO devuelven valores numéricos directos:

\begin{lstlisting}
# INCORRECTO - Causa error de formato
tc = comp.CriticalTemperature
print(f"Tc = {tc:.2f} K")
# Error: unsupported format string passed to CDispatch

# INCORRECTO - Causa error de conversion
tc = float(comp.CriticalTemperature)
# Error: float() argument must be 'string or number',
#        not 'CDispatch'
\end{lstlisting}

\subsubsection{Solución: Usar .GetValue()}

\begin{successbox}
\begin{lstlisting}
# Metodo 1: GetValue() - RECOMENDADO
tc = comp.CriticalTemperature.GetValue()
pc = comp.CriticalPressure.GetValue()
mw = comp.MolecularWeight.GetValue()

print(f"Tc = {tc:.2f} K")    # Funciona
print(f"Pc = {pc:.2f} kPa")
print(f"MW = {mw:.2f} g/mol")

# Metodo 2: Propiedad .Value - ALTERNATIVA
tc = comp.CriticalTemperature.Value  # Tambien funciona
\end{lstlisting}
\end{successbox}

\subsubsection{Exploración de Objeto CDispatch}

Para debugging, explorar la estructura del objeto COM:

\begin{lstlisting}
# Debugging: Explorar propiedades disponibles
valor = comp.CriticalTemperature
print(f"Tipo: {type(valor)}")
# <class 'win32com.client.CDispatch'>
print(f"Valor directo: {valor}")
# 239.44802246093752

# Metodos disponibles:
# - GetValue()  -> float
# - Value       -> float
# - SetValue(val)
# - CalcBy
# - IsKnown
# - UnitConversionType
\end{lstlisting}

\subsection{Manejo de Corrientes (Streams)}

\begin{lstlisting}
flowsheet = case.Flowsheet
streams = flowsheet.MaterialStreams

# Crear corriente
stream1 = streams.Add("Entrada_Metanol")

# Especificar condiciones termodinamicas PRIMERO
stream1.TemperatureValue = 25 + 273.15  # K
stream1.PressureValue = 101.325          # kPa
stream1.MolarFlowValue = 100.0           # kgmole/h

# Especificar composicion usando ComponentMolarFraction.Values
# IMPORTANTE: ComponentMolarFractionValue NO es callable con 2 params
comp_frac = stream1.ComponentMolarFraction.Values
comp_frac[0] = 1.0  # 100% primer componente (Methanol)
stream1.ComponentMolarFraction.Values = comp_frac

# Leer propiedades (requiere GetValue())
T = stream1.TemperatureValue.GetValue() - 273.15
P = stream1.PressureValue.GetValue()
F = stream1.MolarFlowValue.GetValue()
rho = stream1.DensityValue.GetValue()  # Calculada
\end{lstlisting}

\begin{warningbox}
\textbf{ERROR COMUN:} NO usar \texttt{ComponentMolarFractionValue("Methanol", 1.0)}
\\Esto causa error: \texttt{'tuple' object is not callable}
\end{warningbox}

\begin{infobox}
\textbf{Hallazgo:}
\begin{itemize}
    \item \textbf{Composicion:} Usar \texttt{ComponentMolarFraction.Values} (array)
    \item \textbf{Escritura T/P/F:} Asignacion directa sin \texttt{.SetValue()}
    \item \textbf{Lectura:} Requiere \texttt{.GetValue()} para obtener \texttt{float}
    \item \textbf{Indices:} El array sigue el orden de agregacion de componentes
\end{itemize}
\end{infobox}

% Sección 4: Errores COM
\section{Errores COM Comunes y Soluciones}

\subsection{Códigos de Error}

\begin{longtable}{@{}p{3cm}p{4cm}p{4cm}p{3cm}@{}}
\toprule
\textbf{Código COM} & \textbf{Descripción} & \textbf{Causa Común} & \textbf{Solución} \\ \midrule
\endfirsthead

\multicolumn{4}{c}%
{{\bfseries Tabla \thetable\ (continuación)}} \\
\toprule
\textbf{Código COM} & \textbf{Descripción} & \textbf{Causa Común} & \textbf{Solución} \\ \midrule
\endhead

\bottomrule
\endfoot

\bottomrule
\endlastfoot

-2147221021 & Operación no disponible & HYSYS no registrado, arquitectura incompatible & Verificar Python 64-bit, reinstalar HYSYS \\
-2147352567 & Ocurrió una excepción & Nombre de componente inválido, propiedad no definida & Usar "H2O" en lugar de "Water", asignar antes de leer \\
-2147024809 & Argumento inválido & Nombre de PropertyPackage incorrecto & Probar variantes ("SRK", "PR", "NRTL") \\
-2147467259 & Error genérico & Propiedad null/no inicializada & Asignar PropertyPackageName antes de leer \\
\caption{Códigos de error COM comunes}
\label{tab:com_errors}
\end{longtable}

\subsection{Errores de Atributo}

\begin{lstlisting}
# Error: <unknown>.Add
# Causa: Intentar agregar a coleccion de solo lectura
fluid_pkg.Components.Add("Methanol")  # INCORRECTO

# Solucion: Usar ComponentList
comp_list.Components.Add("Methanol")  # CORRECTO
\end{lstlisting}

\begin{lstlisting}
# Error: <unknown>.PropertyPackage
# Causa: Propiedad no existe
fluid_pkg.PropertyPackage = "NRTL"  # INCORRECTO

# Solucion: Usar PropertyPackageName
fluid_pkg.PropertyPackageName = "NRTL"  # CORRECTO
\end{lstlisting}

\begin{lstlisting}
# Error: Property '<unknown>.SaveRequired' cannot be set
# Causa: Propiedad no disponible en automatizacion
case.SaveRequired = False  # INCORRECTO

# Solucion: Simplemente no usar
hysys.Quit()  # CORRECTO
\end{lstlisting}

% Sección 5: Manejo de Rutas
\section{Manejo de Rutas Portátil}

\subsection{Problema: Rutas Hardcodeadas}

\begin{lstlisting}
# INCORRECTO - No portatil entre Windows/Linux
output_file = '/home/user/aspen_python_API/practica2/resultados.json'
\end{lstlisting}

\subsection{Solución: Rutas Relativas con os.path}

\begin{successbox}
\begin{lstlisting}
import os

# Obtener directorio del script actual
script_dir = os.path.dirname(os.path.abspath(__file__))

# Construir ruta de archivo de salida
output_file = os.path.join(script_dir, 'resultados_aspen.json')

# Funciona en Windows y Linux
with open(output_file, 'w') as f:
    json.dump(results, f, indent=2)
\end{lstlisting}
\end{successbox}

% Sección 6: Plantilla Validada
\section{Plantilla de Código Validada}

\subsection{Script Completo Mínimo}

\begin{lstlisting}
#!/usr/bin/env python3
# -*- coding: utf-8 -*-
import win32com.client as win32
import time
import json
import os

def main():
    try:
        # 1. Iniciar HYSYS
        hysys = win32.Dispatch('HYSYS.Application')
        hysys.Visible = True
        case = hysys.SimulationCases.Add()

        # 2. Configurar componentes
        basis_manager = case.BasisManager
        comp_lists = basis_manager.ComponentLists
        comp_list = comp_lists.Add()
        components = comp_list.Components

        components.Add("Methanol")
        components.Add("H2O")

        # 3. Configurar paquete termodinamico
        fluid_pkg = basis_manager.FluidPackages.Add()

        for nombre in ["NRTL", "SRK", "PR"]:
            try:
                fluid_pkg.PropertyPackageName = nombre
                if fluid_pkg.PropertyPackageName == nombre:
                    print(f"Modelo: {nombre}")
                    break
            except:
                continue

        # 4. Crear corriente
        flowsheet = case.Flowsheet
        streams = flowsheet.MaterialStreams
        stream1 = streams.Add("S1")

        stream1.TemperatureValue = 298.15  # K
        stream1.PressureValue = 101.325     # kPa
        stream1.MolarFlowValue = 100.0      # kgmole/h

        # Composicion (usar array, no metodo con 2 params)
        comp_frac = stream1.ComponentMolarFraction.Values
        comp_frac[0] = 1.0  # 100% Methanol
        stream1.ComponentMolarFraction.Values = comp_frac

        time.sleep(2)  # Esperar calculos

        # 5. Extraer resultados
        T = stream1.TemperatureValue.GetValue()
        P = stream1.PressureValue.GetValue()
        rho = stream1.DensityValue.GetValue()

        print(f"T={T-273.15:.2f}C, P={P:.2f}kPa, rho={rho:.2f}kg/m3")

        # 6. Guardar
        script_dir = os.path.dirname(os.path.abspath(__file__))
        output_file = os.path.join(script_dir, 'resultados.json')

        with open(output_file, 'w') as f:
            json.dump({'T_K': T, 'P_kPa': P, 'rho': rho},
                     f, indent=2)

        # 7. Cerrar
        time.sleep(2)
        hysys.Quit()

    except Exception as e:
        print(f"ERROR: {e}")
        import traceback
        traceback.print_exc()

if __name__ == '__main__':
    main()
\end{lstlisting}

% Sección 7: Casos Especiales
\section{Casos Especiales}

\subsection{Componentes con Disponibilidad Limitada}

\begin{lstlisting}
# Intentar agregar componente que puede no estar disponible
try:
    components.Add("Glycerol")
    print("Glycerol agregado")
except Exception as e:
    print(f"Glycerol no disponible: {e}")
    # Continuar sin este componente o usar sustituto
\end{lstlisting}

\subsection{Operaciones con Reactores}

\begin{lstlisting}
# Crear reactor de conversion
operations = flowsheet.Operations
reactor = operations.Add("ConversionReactor")
reactor.Name = "R1"

# Conectar corrientes
reactor.Feeds.Add(feed_stream)
reactor.Outlet = outlet_stream

# Configurar reaccion
rxn_set = reactor.ReactionSet
rxn = rxn_set.Reactions.Add()

# Estequiometria: A + B -> C + D
rxn.ComponentStoichCoeffValue("Methanol", -1.0)   # Reactivo
rxn.ComponentStoichCoeffValue("Ethanol", -1.0)    # Reactivo
rxn.ComponentStoichCoeffValue("Ethyl acetate", 1.0) # Producto
rxn.ComponentStoichCoeffValue("H2O", 1.0)         # Producto

# Conversion fija
rxn.BaseComponent = "Ethanol"
rxn.Conversion = 0.85  # 85%

time.sleep(3)  # Esperar convergencia
\end{lstlisting}

% Sección 8: Mejores Prácticas
\section{Mejores Prácticas}

\subsection{Manejo de Errores}

\begin{lstlisting}
try:
    # Operacion COM
    components.Add("ComponentName")
except Exception as e:
    print(f"Error: {e}")
    print(f"Tipo: {type(e).__name__}")

    # Para errores COM, obtener codigo
    if hasattr(e, 'args') and len(e.args) > 0:
        error_code = e.args[0] if isinstance(e.args[0], int) else None
        print(f"Codigo COM: {error_code}")

    import traceback
    traceback.print_exc()
\end{lstlisting}

\subsection{Esperas para Cálculos}

\begin{infobox}
HYSYS requiere tiempo para convergencia de cálculos:
\begin{itemize}
    \item Después de especificar corrientes: \texttt{time.sleep(2)}
    \item Después de crear reactores: \texttt{time.sleep(3)}
    \item Antes de cerrar HYSYS: \texttt{time.sleep(2)}
\end{itemize}
\end{infobox}

\subsection{Verificación de Conteos}

\begin{lstlisting}
# Siempre verificar que los componentes se agregaron
print(f"Componentes agregados: {components.Count}")

if components.Count == 0:
    raise ValueError("No se agregaron componentes")
\end{lstlisting}

% Sección 9: Debugging
\section{Debugging Sistemático}

\subsection{Explorar Objeto Desconocido}

\begin{lstlisting}
# Listar todas las propiedades/metodos disponibles
obj = basis_manager  # Objeto a explorar

print("Propiedades disponibles:")
for attr in dir(obj):
    if not attr.startswith('_'):
        print(f"  - {attr}")
\end{lstlisting}

\subsection{Verificar Tipo de Valor}

\begin{lstlisting}
valor = comp.CriticalTemperature

print(f"Tipo: {type(valor)}")
print(f"Valor directo: {valor}")

# Si es CDispatch, explorar submetodos
if hasattr(valor, '_oleobj_'):
    print("Es objeto COM")
    for attr in dir(valor):
        if not attr.startswith('_') and attr[0].isupper():
            print(f"  - {attr}")
\end{lstlisting}

\subsection{Script de Diagnóstico}

\begin{lstlisting}
import sys

print("="*60)
print("DIAGNOSTICO ASPEN HYSYS + PYTHON")
print("="*60)

# 1. Arquitectura Python
arch = "64-bit" if sys.maxsize > 2**32 else "32-bit"
print(f"\n1. Python: {sys.version}")
print(f"   Arquitectura: {arch}")

# 2. Libreria pywin32
try:
    import win32com.client
    print("\n2. pywin32 instalado")
except ImportError:
    print("\n2. pywin32 NO instalado")
    print("   Instalar: pip install pywin32")

# 3. Conexion HYSYS
try:
    hysys = win32com.client.Dispatch('HYSYS.Application')
    print("\n3. HYSYS conectado exitosamente")
    hysys.Visible = False
    hysys.Quit()
except Exception as e:
    print(f"\n3. Error conectando HYSYS: {e}")
\end{lstlisting}

% Sección 10: Conclusiones
\section{Conclusiones y Recomendaciones}

\subsection{Hallazgos Clave}

\begin{enumerate}
    \item La documentación estándar de HYSYS \textbf{no refleja la API real} expuesta en todas las versiones
    \item Los nombres de componentes y paquetes termodinámicos \textbf{varían} entre versiones
    \item Los objetos COM requieren \textbf{manejo especial} para extracción de valores numéricos
    \item La jerarquía \texttt{BasisManager → ComponentLists} es \textbf{obligatoria}, no opcional
\end{enumerate}

\subsection{Recomendaciones}

\begin{enumerate}
    \item \textbf{Siempre usar debugging exploratorio} antes de asumir que la documentación es correcta
    \item \textbf{Implementar fallbacks} para nombres de componentes y paquetes termodinámicos
    \item \textbf{Usar \texttt{.GetValue()} consistentemente} para todas las propiedades COM
    \item \textbf{Verificar conteos} después de cada operación de agregación
    \item \textbf{Incluir tiempos de espera} para permitir convergencia de cálculos
\end{enumerate}

\subsection{Recursos de Referencia}

\begin{itemize}
    \item \textbf{Código fuente validado:} Prácticas 2-8 en repositorio
    \item \textbf{Script de diagnóstico:} \texttt{practica1/diagnostico\_hysys.py}
    \item \textbf{Guía de troubleshooting:} \texttt{practica1/SOLUCION\_PROBLEMAS.md}
\end{itemize}

% Apéndice
\appendix

\section{Resumen de Commits}

Durante la sesión de debugging se realizaron los siguientes commits:

\begin{itemize}
    \item \texttt{f7c2922} - Corregir rutas hardcodeadas en todos los archivos py\_alone.py
    \item \texttt{d4bc49b} - Aplicar patrones correctos de HYSYS API a prácticas 6-8
    \item \texttt{3cf4d61} - Aplicar patrones correctos de HYSYS API a prácticas 2-5
    \item \texttt{8b3077c} - Usar GetValue() para extraer valores numéricos de propiedades COM
    \item \texttt{23b1881} - Agregar debug para explorar estructura de propiedades de componentes COM
\end{itemize}

\section{Glosario}

\begin{description}
    \item[COM] Component Object Model - Tecnología de Microsoft para comunicación entre procesos
    \item[CDispatch] Objeto COM dinámico en Python (pywin32)
    \item[BasisManager] Objeto de HYSYS que gestiona componentes y paquetes termodinámicos
    \item[FluidPackage] Configuración de modelo termodinámico en HYSYS
    \item[ComponentList] Colección de componentes químicos en HYSYS
    \item[pywin32] Librería Python para acceder a APIs COM de Windows
\end{description}

% Información de versión
\vfill
\hrule
\vspace{0.5cm}
\noindent
\textbf{Versión del documento:} 1.0 \\
\textbf{Última actualización:} 2025-01-20 \\
\textbf{Mantenedor:} Claude (Anthropic) \\
\textbf{Estado:} Validado mediante 8 prácticas funcionales \\
\textbf{Repositorio:} \texttt{aspen\_python\_API} \\
\textbf{Rama:} \texttt{claude/debug-practice-1-01MPgRrgcABHGA1GhADpKGPJ}

\end{document}
